% !TEX program = pdflatex
% Cogniland project overview: environment, methods, experiments, results.
% Self-contained, stock packages only.
\documentclass[11pt]{article}

\usepackage[margin=1in]{geometry}
\usepackage{graphicx}
\usepackage{booktabs}
\usepackage{amsmath}
\usepackage{amssymb}
\usepackage{xcolor}
\usepackage{caption}
\usepackage{float}
\usepackage{enumitem}
\usepackage[colorlinks=true,linkcolor=violet,citecolor=teal,urlcolor=teal]{hyperref}
\usepackage[section]{placeins}

\definecolor{lakes}{RGB}{42,120,214}
\definecolor{rocky}{RGB}{217,89,38}
\definecolor{balanced}{RGB}{25,158,112}
\captionsetup{font=small, labelfont=bf}
\newcommand{\code}[1]{\texttt{\small #1}}
\newcommand{\lakes}{\textcolor{lakes}{\textbf{lakes}}}
\newcommand{\rocky}{\textcolor{rocky}{\textbf{rocky}}}
\newcommand{\balanced}{\textcolor{balanced}{\textbf{balanced}}}
\newcommand{\R}{\mathbb{R}}
\newcommand{\E}{\mathbb{E}}
\newcommand{\ctg}{\mathrm{ctg}}

\title{\textbf{Cogniland}\\[3pt]
\large A Partially-Observed Navigation Testbed for the\\
Mechanistic Study of \emph{Belief} in World Models and Recurrent Agents}
\author{Project overview}
\date{\today}

\begin{document}
\maketitle
\vspace{-2em}

\begin{abstract}\noindent
\textbf{Cogniland} is a research codebase built around one idea: give a
reinforcement-learning agent a task where it \emph{cannot see everything at
once}, so that to act well it must build and hold an internal \emph{belief}
about the hidden state of the world; then open the agent up and study that
belief directly. This document is a self-contained tour of the project. We
describe the current flagship environment---the \code{bridge\_tunnel}
\code{fork\_wall} task, a small grid world in which a hidden ``map category''
must be inferred from terrain seen early, remembered across a blind corridor,
and finally read out as a left-or-right door choice. We then describe the two
kinds of agents we train on it (a recurrent policy and a \emph{world model}),
the interpretability tools we use to read and edit their beliefs (linear
probes, activation steering, and manifold geometry), and the experiments and
results those tools produce. Throughout we favour plain language and give the
math needed to reproduce every claim. The recurring lesson is that the agent's
belief is real, causal, and geometrically structured---but also that getting a
world model to learn the task at all exposed a subtle reward-design pitfall
that we document and fix.
\end{abstract}

\tableofcontents
\vspace{1em}

% ======================================================================
\section{Introduction: why a ``belief'' testbed}
\label{sec:intro}

Most interesting decision problems are \emph{partially observed}: the thing you
need to know is not written on the current observation. A driver at a fork does
not see the destination; a rat in a maze does not see the reward arm; a language
model predicting the next token does not see the author's intent. Formally these
are \emph{partially observable Markov decision processes} (POMDPs)
\cite{kaelbling1998planning}. The classical result for POMDPs is that the
optimal action depends not on the current observation $o_t$ but on the
\emph{belief state}---a probability distribution over the hidden state given the
entire history $o_{1:t}, a_{1:t-1}$. A good agent must therefore compute
something belief-like and carry it forward in memory.

Recurrent neural networks do exactly this in practice: their hidden state is a
learned, compressed summary of history, and there is now strong evidence---in
neuroscience \cite{mante2013context,wong2006recurrent,sussillo2013opening}, in
meta-RL \cite{mikulik2020meta,lambrechts2022recurrent}, and in large models
\cite{li2023othello,shai2024transformers}---that these hidden states encode
recoverable, sometimes even Bayes-optimal, beliefs about latent variables.
\textbf{Cogniland exists to make this phenomenon concrete and manipulable.} We
want a task where (i) the hidden variable is simple and known to us (so we can
check what the agent believes against the truth), (ii) the belief must genuinely
be \emph{held in memory} (so it is not a trivial reflex), and (iii) the belief
must be \emph{used} to drive a discrete decision (so it is behaviourally
load-bearing). The \code{fork\_wall} environment (\S\ref{sec:env}) is engineered
to satisfy all three.

Once such a task exists, we can ask sharper questions than ``does the agent
succeed?''. We can ask: \emph{where} is the belief stored? Is it \emph{causal}
(does editing it change behaviour)? Does it lie on a simple geometric object?
And in a \emph{world model}---an agent that learns to imagine the future---does
the belief shape what the agent dreams? Sections \ref{sec:methods} and
\ref{sec:experiments} develop the tools and answers.

% ======================================================================
\section{The environment: \code{bridge\_tunnel} / \code{fork\_wall}}
\label{sec:env}

\subsection{The task in one paragraph}
The agent is a single point that starts at the left edge of a rectangular grid
and must reach a goal door on the right. Early in its journey it walks past
patches of \textbf{terrain}---water and rock---whose \emph{amounts} secretly
encode the map's \emph{category}. It then crosses a stretch of empty grass (the
\emph{memory corridor}), where no terrain is visible, and arrives at a
\textbf{wall} with a narrow passage. Beyond the wall are two doors, one near the
top and one near the bottom. Exactly which door counts as success depends on the
category the agent saw earlier. Because the terrain has scrolled out of view by
the time the agent reaches the doors, it must have \emph{remembered} the
category to choose correctly. Figure~\ref{fig:maps} shows one map per category.

\begin{figure}[H]
\centering
\includegraphics[width=\textwidth]{../outputs/bridge_tunnel_forkwall/sample_maps.png}
\caption{Three maps from the fixed training set, one per category. Spawn is the
white marker on the left. Category evidence (water for \lakes, rock for \rocky,
both for \balanced) is confined to the left of the cyan line; a grass
\emph{memory corridor} then runs to the fork wall and its top/bottom doors
(yellow). By the fork the evidence has left the agent's egocentric view, so the
door choice requires memory rather than a glance.}
\label{fig:maps}
\end{figure}

\subsection{The hidden variable: map category}
Every map belongs to one of three categories, defined by how much water versus
rock its terrain contains:
\begin{itemize}[nosep]
\item \lakes{}: water-dominant. The \emph{correct} door is the \textbf{bottom}.
\item \rocky{}: rock-dominant. The correct door is the \textbf{top}.
\item \balanced{}: roughly equal water and rock. \emph{Either} door counts.
\end{itemize}
The category is never handed to the agent as an input; it is only implicit in
the terrain. It is the ``hidden state'' of our POMDP, and ``what the agent
believes the category is'' is the belief we study throughout.

\subsection{Observation: an egocentric window}
The agent never sees the whole map. At each step it receives a square,
agent-centred crop of the terrain of side $V=21$ cells, together with a few
scalar signals (which way the goal lies, how far into the episode we are). Write
the terrain as $T\in\{0,\dots,K-1\}^{H\times W}$ with $K=9$ tile types (grass,
water, rock, wood, target, tree, sand, dirt, and an out-of-bounds pad), $H=32$,
$W=64$. The observation is
\begin{equation}
o_t \;=\; \big(\underbrace{\mathrm{onehot}\big(T[\,r_t-\tfrac{V-1}{2}:r_t+\tfrac{V-1}{2}\,,\ c_t-\tfrac{V-1}{2}:c_t+\tfrac{V-1}{2}]\big)}_{V\times V\times K \text{ one-hot crop}},\ \underbrace{s_t}_{\in\R^{5}}\big),
\end{equation}
where $(r_t,c_t)$ is the agent's cell and cells outside the map are filled with
the out-of-bounds tile. Flattened, this is a vector of dimension
$V^2K + 5 = 21^2\cdot 9 + 5 = 3974$. Crucially, because the window has radius
$\lfloor V/2\rfloor = 10$, terrain more than $10$ cells behind the agent is
\emph{invisible}. This single fact is what makes the memory corridor
(\S\ref{sec:memory}) do its job.

\subsection{Actions and dynamics}
The action space is $\mathcal{A}=\{\text{up},\text{down},\text{left},
\text{right},\text{build},\text{mine}\}$ (six discrete actions). Moves are
grid steps; \code{build} turns a faced water cell into a walkable bridge and
\code{mine} turns a faced rock cell into grass. In the variant used here
(``\code{BT}-rules'', i.e.\ commitment disabled) the tools are always available
and the agent can also simply walk around obstacles, so the decision that
matters is purely \emph{which door} it goes to. An episode ends when the agent
steps onto either door, or after a step budget of $800$.

\subsection{The memory corridor (task design)}
\label{sec:memory}
An earlier version of this map placed terrain right up to the fork, so the agent
could see the evidence and the doors \emph{in the same window} and pick the door
by reflex---no memory needed. We fix this with a \textbf{memory corridor}: a band
of $g=16$ pure-grass columns immediately before the wall, with all category
terrain confined to its left. Let $c_{\text{wall}}$ be the wall column and
$c_{\text{ev}}$ the rightmost terrain column; the construction guarantees
$c_{\text{wall}}-c_{\text{ev}} \ge g$. Since the view radius is $10 < g$, when
the agent stands at the fork ($c_t \approx c_{\text{wall}}-1$) the nearest
evidence at $c_{\text{ev}}$ is $\ge g - 1 > 10$ cells away and therefore outside
the window. Concretely the evidence leaves view once $c_t - 10 > c_{\text{ev}}$,
giving the agent on the order of $c_{\text{wall}} - (c_{\text{ev}}+10)\approx 6$
``blind'' timesteps in which the correct action depends only on memory. This
turns a reflex into a genuine working-memory problem.

\subsection{Reward, and a reward-shaping pitfall}
\label{sec:reward}
The reward at each step is a sum of three parts:
\begin{equation}
r_t \;=\; \underbrace{-\lambda_{\text{slack}}}_{\text{time cost}}
\;+\; \underbrace{\rho\,\big[\text{stepped onto the correct door}\big]}_{\text{goal bonus}}
\;+\; \underbrace{\kappa\,\big(\Phi(s_{t})-\gamma_{\!\text{sh}}\,\Phi(s_{t+1})\big)}_{\text{shaping}},
\label{eq:reward}
\end{equation}
with $\lambda_{\text{slack}}=0.01$, $\rho=3.0$, $\kappa=0.015$, and the
\emph{potential} $\Phi(s) = -\ctg(s)$, where $\ctg(s)$ is the number of steps
from the agent's cell to the correct door (a cost-to-go computed by
breadth-first search on the terrain). Shaping of the form
$F = \gamma_{\!\text{sh}}\Phi(s') - \Phi(s)$ is the classic
\emph{potential-based reward shaping} of Ng, Harada, and Russell
\cite{ng1999policy}, which provably does not change the optimal policy when
$\gamma_{\!\text{sh}}$ equals the RL discount $\gamma$.

\paragraph{The pitfall.} Potential-based shaping is policy-invariant in theory
but not \emph{exploration}-invariant in practice. At a fixed state
($s_{t+1}=s_t$, i.e.\ the agent stands still), the shaping term is
\begin{equation}
\kappa\big(\Phi - \gamma_{\!\text{sh}}\Phi\big)
= \kappa(1-\gamma_{\!\text{sh}})\,\big(-\Phi\big)
= \kappa(1-\gamma_{\!\text{sh}})\,\ctg(s) \;>\; 0 .
\end{equation}
Near the spawn, $\ctg$ is large ($\approx 70$), so with $\gamma_{\!\text{sh}}
=\gamma=0.99$ this is $0.015\cdot0.01\cdot70 \approx 0.0105$ per step, which
\emph{exceeds} the time cost $\lambda_{\text{slack}}=0.01$. \textbf{The agent is
literally paid to do nothing.} We observed exactly this failure: a DreamerV3
world model, trained from scratch, converged to a policy that never left the
spawn tile (maximum column reached $=0$) while collecting positive reward and
timing out every episode. On-policy PPO happened to explore its way out of this
basin, but the world model, which learns its policy \emph{in imagination}, sat
in it comfortably; raising the exploration bonus did not help.

\paragraph{The fix.} We set $\gamma_{\!\text{sh}}=1$, turning the shaping into
\emph{pure progress}:
\begin{equation}
\kappa\big(\Phi(s_t)-\Phi(s_{t+1})\big)
= \kappa\big(\ctg(s_t)-\ctg(s_{t+1})\big).
\end{equation}
Now standing still earns exactly zero shaping (net reward $-\lambda_{\text{slack}}
<0$), a step toward the goal earns $+\kappa$, and a step away earns $-\kappa$;
there is no stationary way to accumulate reward, so the only route to positive
return is to reach a door. Empirically this immediately removed the basin: the
world model began leaving the spawn and reaching doors, and evaluation success
rose from a flat $0$ toward the same regime PPO reaches. We flag this because it
is a general and easily-missed lesson: \emph{discounted potential shaping can
create a ``loitering'' optimum whenever the potential is large and the step cost
is small}, and an imagination-based learner is especially prone to fall into it.

\subsection{A fixed map dataset}
So that every agent trains and is tested on identical data, we pre-generate a
fixed dataset of $6000$ maps---$2000$ per category---and split it $80/20$ into
$4800$ \emph{train} and $1200$ \emph{test} maps. Maps are procedurally generated
from simplex-noise heightmaps thresholded to hit the target water/rock fractions,
then post-processed to carve the memory corridor and guarantee the correct door
is reachable. Training draws maps uniformly from the train split; all reported
success numbers are on the disjoint test split, so nothing evaluated was ever
seen in training.

% ======================================================================
\section{The agents}
\label{sec:agents}

We study two agents on identical data. They differ in \emph{what} they learn:
the first learns only to act; the second additionally learns to \emph{imagine}.

\subsection{Recurrent PPO with a GRU memory}
The first agent is a recurrent policy trained with Proximal Policy Optimization
(PPO) \cite{schulman2017proximal}. Because the task is partially observed, the
policy needs memory; we give it a Gated Recurrent Unit (GRU) \cite{cho2014gru}.
An encoder maps the observation $o_t$ to a feature $e_t$; the GRU updates a
hidden state
\begin{equation}
h_t = \mathrm{GRU}(e_t, h_{t-1}) \in \R^{128},
\end{equation}
and linear heads read the action distribution $\pi_\theta(a_t\mid h_t)$ and value
$V_\theta(h_t)$ off $h_t$. The GRU hidden state $h_t$ \emph{is} the agent's
working memory and hence its belief; it is what a probe (\S\ref{sec:probe})
reads. PPO optimizes the clipped surrogate
\begin{equation}
\mathcal{L}^{\text{PPO}}(\theta)=\E_t\Big[\min\big(\varrho_t(\theta)\hat A_t,\
\mathrm{clip}(\varrho_t(\theta),1-\epsilon,1+\epsilon)\hat A_t\big)\Big],\quad
\varrho_t(\theta)=\frac{\pi_\theta(a_t\mid h_t)}{\pi_{\theta_{\text{old}}}(a_t\mid h_t)},
\end{equation}
with advantages $\hat A_t$ from generalized advantage estimation. We optionally
add a small \emph{auxiliary head} that predicts the map category from $h_t$; this
does not change behaviour but makes the belief more linearly readable, and its
accuracy (\code{belief\_acc}) is a convenient training-time diagnostic.

\subsection{World models: DreamerV3 / R2-Dreamer}
The second agent is a \emph{world model} in the DreamerV3 family
\cite{hafner2023dreamerv3,ha2018world}. A world model learns a compact latent
model of the environment's dynamics and then trains its policy by ``dreaming''
---rolling that latent model forward and doing RL on the imagined trajectories,
never touching the real environment during policy improvement. Concretely we use
the in-repo \code{purejaxwm} port and an efficient PyTorch reimplementation,
R2-Dreamer \cite{morihira2026r2dreamer}. The core is a \emph{Recurrent State
Space Model} (RSSM), whose state has a deterministic part $d_t$ and a stochastic
part $z_t$:
\begin{align}
d_t &= f_\phi(d_{t-1}, z_{t-1}, a_{t-1}) && \text{(deterministic recurrence, a GRU)}\\
\hat z_t &\sim p_\phi(\hat z_t \mid d_t) && \text{(\emph{prior}: what the model expects, no observation)}\\
z_t &\sim q_\phi(z_t \mid d_t, e_t) && \text{(\emph{posterior}: corrected by the encoded observation $e_t$)}
\end{align}
The world model is trained to (a) predict the next stochastic state (a KL
between posterior and prior), (b) predict reward and episode continuation, and,
in the reconstruction variant we use here, (c) \emph{decode} the observation back
from the latent state, $\hat o_t = \mathrm{dec}_\phi(z_t, d_t)$. Writing the full
latent (``feature'') as $b_t = [\,z_t;\,d_t\,]\in\R^{3840}$, the world-model loss
is a variational bound of the form
\begin{equation}
\mathcal{L}^{\text{WM}} = \E\Big[\underbrace{-\log \mathrm{dec}_\phi(o_t\mid b_t)}_{\text{reconstruction}}
\;+\;\underbrace{-\log p_\phi(r_t\mid b_t)-\log p_\phi(\text{cont}_t\mid b_t)}_{\text{reward/continue}}
\;+\;\beta\,\underbrace{\mathrm{KL}\!\big[q_\phi(z_t)\,\|\,p_\phi(\hat z_t)\big]}_{\text{dynamics}}\Big].
\end{equation}
The actor and critic are trained purely on latent rollouts of the \emph{prior}
(imagination): starting from a real $b_t$, one repeatedly samples
$a\sim\pi$, advances $d,\hat z$ with $f_\phi$ and $p_\phi$, and optimizes the
policy against imagined reward. This is what makes a world model special for us:
its belief $b_t$ is not only read by the policy---it \emph{generates} the future
the policy plans against. We can literally decode what the agent imagines.

\paragraph{Reconstruction vs.\ decoder-free.} R2-Dreamer's default objective is
decoder-free (it replaces reconstruction with a redundancy-reduction term). For
this project we deliberately train with the \emph{reconstruction} objective
(``\code{rep\_loss=dreamer}'') so that the world model is genuinely generative
and its dreams decode to legible maps---at the mild cost of extra compute. When
we analyse an older decoder-free checkpoint, we instead fit a small decoder
post hoc by frozen-backbone probing (it reaches $99.6\%$ per-cell tile accuracy),
which recovers the same picture.

\subsection{What we mean by ``belief''}
Two notions of belief appear, and it helps to keep them straight.
\begin{itemize}[nosep]
\item The \textbf{internal belief} is the recurrent state itself---the GRU hidden
$h_t$ for PPO, the RSSM feature $b_t=[z_t;d_t]$ for the world model. It is a
high-dimensional vector, the physical substrate of memory.
\item The \textbf{belief distribution} is a probability vector over the three map
categories, $\mathbf{p}_t = (p_t^{\balanced}, p_t^{\lakes}, p_t^{\rocky})\in
\Delta^2$. This is the ``Bayesian'' belief a reader expects: as evidence
accumulates it should sharpen from near-uniform (uncertain) toward a one-hot
(committed). We obtain it by \emph{probing} the internal belief
(\S\ref{sec:probe}).
\end{itemize}
The whole project is, in a sense, about the map from the first to the second, and
about whether that map is causal and geometric.

% ======================================================================
\section{Methods: reading and editing belief}
\label{sec:methods}

\subsection{Probing: decoding the belief distribution}
\label{sec:probe}
A \emph{probe} is a simple classifier fit to predict the hidden category from the
internal belief \cite{ravichander2021probing,li2023othello}. Given a dataset of
(belief, true-category) pairs $\{(b^{(i)}, c^{(i)})\}$ collected from rollouts,
we fit two linear probes:
\begin{enumerate}[nosep]
\item \textbf{Nearest class mean.} Compute the per-category centroid
$\mu_c=\frac{1}{n_c}\sum_{i:c^{(i)}=c} b^{(i)}$ and predict
$\hat c(b)=\arg\min_c \|b-\mu_c\|_2$. This is the simplest possible linear
readout; its decision boundary is the perpendicular bisector between centroids.
\item \textbf{Logistic regression.} Fit a multinomial logistic model
$\mathbf{p}(b)=\mathrm{softmax}(Wb+w_0)$ with an $\ell_2$ penalty chosen by
cross-validation. This is a \emph{discriminative} linear probe: it learns the
directions that best separate the classes, and its softmax output is exactly the
belief distribution $\mathbf{p}_t$.
\end{enumerate}
The probe is deliberately linear: if a linear map recovers the category, the
belief is stored in a linearly-accessible way, which is the property that makes
the later steering experiments meaningful. Because the probe is trained on one
set of episodes and evaluated on held-out ones, its accuracy measures genuine
\emph{decodability}, not memorization.

\subsection{Steering: editing the belief to test causality}
\label{sec:steer}
Decoding shows the belief is \emph{present}; steering tests whether it is
\emph{causal}. The idea, following activation-editing work in language models
\cite{turner2023activation,zou2023representation,meng2022locating} and the
causal-abstraction programme \cite{geiger2021causal}, is to \emph{patch} the
belief toward another category and see whether behaviour follows. We form a
steering direction $\Delta_{s\to t}$ from source category $s$ to target $t$ and
apply
\begin{equation}
b' \;=\; b \;+\; \alpha\,\Delta_{s\to t},
\end{equation}
then feed $b'$ to the (frozen) actor and record its decision. Two choices of
direction:
\begin{itemize}[nosep]
\item \textbf{Difference of means:} $\Delta_{s\to t}=\mu_t-\mu_s$, the vector
between class centroids.
\item \textbf{Probe weights:} $\Delta_{s\to t}=w_t-w_s$, the difference of the
logistic probe's per-class weight vectors---the directions the classifier
actually found most discriminative.
\end{itemize}
A clean causal result looks like: patching a belief from \rocky\ toward \lakes\
flips the agent's door choice from top to bottom. Reporting both directions lets
us separate ``the belief is decodable'' (a property of the probe) from ``the
belief is steerable'' (a property of the representation)---these are not the
same, as \S\ref{sec:exp1} shows.

\subsection{Geometry: is belief a direction or a manifold?}
\label{sec:manifold}
Difference-of-means steering assumes the belief varies along a straight
\emph{line} between categories. But recurrent belief states are known to trace
\emph{curved, low-dimensional structures}---line attractors and rings---in their
activation space \cite{sussillo2013opening,shai2024transformers}. If the belief
lives on a curved \emph{manifold}, moving in a straight line between two beliefs
leaves the manifold, landing on states the network never actually occupies. We
test this directly, adapting the manifold-steering method of the Goodfire
\code{causalab} framework \cite{geiger2021causal}.

Concretely we parameterize belief by a scalar \emph{evidence coordinate}
$\zeta\in[-1,1]$ (signed by which door the trajectory heads to, scaled by
progress), bin it, and take the mean belief in each bin as a control point
$\{(\zeta_j, \bar b_j)\}$. We fit a natural cubic smoothing spline
$\gamma:[-1,1]\to\R^{3840}$ through these points (the 1-D analogue of a
thin-plate spline, minimizing curvature $\int\|\gamma''\|^2$)
\cite{reinsch1967smoothing}, giving a smooth 1-D curve $\mathcal{M}=\{\gamma(u)\}$
in belief space. Two ways to steer from belief $\gamma(u_a)$ to $\gamma(u_b)$:
\begin{align}
\text{linear:}\quad & b(t) = (1-t)\,\gamma(u_a) + t\,\gamma(u_b), \\
\text{manifold:}\quad & b(t) = \gamma\big((1-t)u_a + t\,u_b\big),
\end{align}
i.e.\ interpolate in \emph{ambient space} (a chord) versus in the \emph{intrinsic
coordinate} (staying on the curve). We measure how far each path strays from the
manifold via Gauss--Newton projection: for a point $x$, find
$u^\star=\arg\min_u\|\gamma(u)-x\|^2$ and report the residual
$\|x-\gamma(u^\star)\|$. If the belief manifold is curved, the linear chord will
bow off it while the manifold path stays on it by construction.

\subsection{Dreams: decoding the imagined future}
\label{sec:dreams}
Because the world model is generative, we can ask what it \emph{imagines}. From a
belief $b_t$ we roll the prior forward under the actor for $H$ steps and decode
each latent to a map crop:
\begin{equation}
\hat o_{t+k} = \mathrm{dec}_\phi(b_{t+k}), \qquad
b_{t+k+1} = \big[\,\hat z_{t+k+1};\,f_\phi(d_{t+k}, \hat z_{t+k}, a_{t+k})\,\big],\quad a_{t+k}\sim\pi.
\end{equation}
This yields a sequence of imagined observations---a ``dream''---that we can look
at directly and measure (e.g.\ what fraction of imagined tiles are water vs.\
rock, or where the goal door appears). Dreams let us ask the sharpest version of
our question: does the belief not only label the map but \emph{generate} a future
consistent with that label?

% ======================================================================
\section{Experiments and results}
\label{sec:experiments}

We first establish that the task is solvable and that belief is the mechanism,
then progressively test how structured and how causal that belief is. Unless
noted, analyses are run on a trained \code{fork\_wall} \code{BT}-rules agent and
report held-out numbers.
The figures in this section are produced by the belief-analysis pipeline on a
trained \code{fork\_wall} \code{BT}-rules agent (held-out episodes); the memory-corridor
world model of \S\ref{sec:reward} is retrained on the fixed dataset and its
figures will be regenerated identically once it saturates.
\newcommand{\cogfig}{figures/cogniland}

\subsection{Training: the task is solvable, and requires memory}
On the fixed dataset with the memory corridor, recurrent PPO reaches
\textbf{$100\%$ success} on held-out maps, with an auxiliary category-probe
accuracy of $\approx 0.90$ and efficient $\sim\!100$-step episodes---it walks the
corridor, holds the category across the blind stretch, and picks the correct
door. That the corridor makes this a genuine memory problem is visible in
training: on the old evidence-at-the-fork maps the policy hit $100\%$ almost
immediately, whereas on the memory maps it climbs gradually as the recurrent
state learns to carry the belief. The world model, once the reward-shaping
pitfall of \S\ref{sec:reward} is fixed, likewise leaves the spawn and its success
rises from $0$ into the same regime. The key methodological takeaways are (i) the
environment is well-posed and solvable, and (ii) an imagination-based learner is
uniquely sensitive to farmable shaping and must be trained with the
pure-progress form.

\subsection{Experiment 1: the belief \emph{causes} the fork decision}
\label{sec:exp1}
We collect beliefs at the moment just before the agent commits up or down, fit
the probes of \S\ref{sec:probe}, and test steering.

\paragraph{The belief is decodable.} The map category is linearly readable from
the pre-decision belief well above the $33\%$ chance level: the nearest-mean
probe reaches $\approx 60\%$ and the logistic probe $\approx 67\%$. (These are
deliberately weak linear readouts; a small nonlinear probe would score higher.)

\paragraph{Behaviour follows the \emph{belief}, not the map.} The decisive test
is a within-model consistency check. On held-out episodes we compare the door the
agent actually reached against (a) the true category and (b) the probe's
\emph{decoded} category. The door matched the \emph{decoded} belief
$\approx\!100\%$ of the time---slightly \emph{tighter} than it matched the ground
truth ($\approx\!92\%$). In other words, on the episodes where the model's belief
was wrong, its \emph{decision} was wrong in exactly the same way. The agent acts
on what it believes, not on what is true---the signature of a genuine internal
belief driving behaviour.

\paragraph{Decodable is not the same as steerable.} Patching the pre-decision
belief with the difference-of-means direction changes the action $\approx\!92\%$
of the time but only lands on the \emph{target} door $\approx\!12\%$ of the time;
using the logistic probe's own discriminative direction raises the on-target rate
to $\approx\!38\%$ (a $3\times$ improvement). Two lessons: the category direction
found by a discriminative probe is a much better steering handle than the naive
centroid difference; and even so, the belief is not perfectly separable by linear
arithmetic---category information is entangled with the ``where am I in the
episode'' information the recurrent state also carries.

\begin{figure}[H]
\centering
\includegraphics[width=\textwidth]{\cogfig/probes_success.png}
\caption{Experiment~1, quantitative. \textbf{(a)} Held-out success per category
for the analysed agent. \textbf{(b)} The map category is linearly decodable from
the pre-decision belief, well above the $33\%$ chance line; a discriminative
(logistic) probe beats the nearest-mean readout. \textbf{(c)} Behaviour follows
the belief: the door reached matches the probe's \emph{decoded} category
($\approx\!100\%$) even more tightly than the true category ($\approx\!92\%$).
Steering by the logistic probe's own direction lands on the target door about
$3\times$ more often than the naive centroid difference---but neither is perfect,
confirming decodable $\ne$ steerable.}
\label{fig:exp1}
\end{figure}

\begin{figure}[H]
\centering
\includegraphics[width=0.86\textwidth]{\cogfig/confusion_door.png}
\caption{Experiment~1, structure. \textbf{(a)} Confusion matrix of the logistic
probe (counts): the diagonal dominates, with the residual confusion concentrated
on the ambiguous \balanced\ class. \textbf{(b)} The category$\to$door behaviour
the belief drives: \rocky\ goes to the top door, \lakes\ to the bottom, \balanced\
to either; timeouts are rare.}
\label{fig:confdoor}
\end{figure}

\begin{figure}[H]
\centering
\includegraphics[width=\textwidth]{\cogfig/belief_pca_end3d.png}
\caption{The belief state \emph{at the end of the episode} (the pre-decision
state at the fork), projected onto its first three principal components.
\textbf{(a)} In 3-D the three categories occupy distinct regions---\rocky\
(stars mark class means) separates cleanly from \lakes/\balanced\ along PC1---so
the category is present in the belief right where the door choice is made.
\textbf{(b)} The same points in the PC1--PC2 plane, now with the marker showing
the \emph{door reached}: up-triangles (top) and down-triangles (bottom) segregate
by category, and the rare timeouts ($\times$) cluster together---the geometry of
the belief and the geometry of the decision line up. Note the top three PCs hold
only $\sim\!30\%$ of the variance: the belief is genuinely high-dimensional, with
category a modest but well-separated factor.}
\label{fig:pca_end}
\end{figure}

\subsection{Experiment 2: the belief shapes the dream}
\label{sec:exp2}
We seed the world model with beliefs from mid-corridor and let it imagine forward
(\S\ref{sec:dreams}), then measure the terrain it hallucinates. Imagined rollouts
from a \lakes\ belief contain markedly more water than those from any other
belief, and rollouts from a \rocky\ belief contain markedly more rock---the
imagined future reflects the believed category, not the currently-visible tiles
(which by mid-corridor are mostly empty grass). The belief is therefore not just
a label read at decision time; it is an active generator of the content of the
agent's planning rollouts.

\begin{figure}[H]
\centering
\includegraphics[width=0.52\textwidth]{\cogfig/dream_composition.png}\hfill
\includegraphics[width=0.46\textwidth]{\cogfig/offmanifold.png}
\caption{Left (Experiment~2): the terrain the world model \emph{imagines} depends
on the believed category---a \lakes\ belief dreams more water, a \rocky\ belief
more rock (mean $\pm$ s.d.\ over imagined rollouts). Right (Experiment~3): a
straight (linear) steering path bows off the fitted belief manifold at
mid-traverse, while the geodesic path stays on it by construction---direct
evidence that the belief is curved.}
\label{fig:dream_off}
\end{figure}

\begin{figure}[H]
\centering
\includegraphics[width=\textwidth]{\cogfig/dream_rollout.png}
\caption{Experiment~2, qualitative. Decoded imagined-future (``dream'') strips,
one row per belief category, seeded mid-corridor and rolled forward under the
actor. The frames are legible maps whose composition matches the seeded belief.}
\label{fig:rollout}
\end{figure}

\subsection{Experiment 3: belief is a \emph{curved} manifold, but a robust one}
\label{sec:exp3}
Using the machinery of \S\ref{sec:manifold} we ask whether belief is a direction
or a manifold. The answer to the geometric question is clear: it is
\textbf{curved}. A straight (linear) steering path between two decisive beliefs
bows measurably off the fitted belief manifold---peaking around $0.5$ in the
evidence-accumulation corridor and up to $\approx\!0.9$ across the full
door-to-door traverse---while the manifold-trace path stays at $0$ off the
manifold by construction. Since a straight line is exactly what direction-based
steering assumes, this explains why that steering was imperfect in
Experiment~1: the direction is a chord across a curved object.

\begin{figure}[H]
\centering
\includegraphics[width=0.82\textwidth]{\cogfig/belief_pca.png}
\caption{Experiment~3: PCA of the belief. Each dot is a real belief state,
coloured by the evidence coordinate $\zeta$ (rocky $-1$ $\to$ neutral $0$ $\to$
lakes $+1$); PC1 of the manifold's own plane tracks $\zeta$ ($r=0.73$). The
fitted spline (black) is the 1-D belief manifold. Steering from the rocky end to
the lakes end along the \emph{geodesic} (green) follows the curve, whereas the
\emph{linear} chord (dashed) cuts across it and bows off---the geometric root of
why direction-based steering is imperfect.}
\label{fig:pca}
\end{figure}

\begin{figure}[H]
\centering
\includegraphics[width=0.66\textwidth]{\cogfig/manifold_3d.png}
\caption{The same corridor belief manifold in three dimensions (a fresh 3-D PCA
of the corridor belief cloud, not the planar curve basis of
Fig.~\ref{fig:pca}). The binned-mean polyline (black) traces a smooth curve
through the cloud from the rocky pole to the lakes pole; the third component
carries little additional structure, confirming the manifold is essentially
one-dimensional, embedded in a high-dimensional, position-dominated space.}
\label{fig:manifold3d}
\end{figure}

However, the striking downstream consequence found for a \emph{pixel} world
model---that leaving the activation manifold corrupts the decoded observation
into an incoherent superposition---does \emph{not} transfer here. On both the
linear and manifold paths the decoded map stays crisp, and, when we steer the
actual fork decision and imagine forward, the predicted door flips cleanly from
one side to the other (the imagined commitment
$p(\text{down})-p(\text{up})$ sweeps from strongly one-sided to the other,
crossing zero) with \emph{no visible corruption}, even where the linear path sits
far off the manifold. Two model-specific reasons: this RSSM belief manifold is
high-dimensional and diffuse (the category coordinate is a low-variance factor
inside a much larger position-dominated state, not the tidy low-dimensional
closed curve a vision encoder traces), and the symbolic decoder and dynamics
generalize gracefully to off-manifold beliefs rather than smearing. The honest
summary is: \emph{the belief geometry is genuinely curved, but robust}---manifold
awareness confirms the shape without being required to steer it in this model.

\subsection{Reading the whole picture}
Putting the experiments together yields a coherent account of the agent's mind on
this task. Terrain evidence early in the corridor moves the internal belief along
a curved, largely one-dimensional path parameterized by ``how much the evidence
favours each door''; a linear probe reads a sharpening category distribution off
that path; the belief is carried, essentially unchanged, across the blind memory
corridor; at the fork the actor reads the belief and commits to the matching
door; and, in the world model, the same belief generates an imagined future whose
terrain and door placement match the believed category. Editing the belief edits
the decision (imperfectly, because category and episode-phase are entangled), and
the belief lives on a curved but robust manifold rather than a straight line.

% ======================================================================
\section{Discussion and lessons}
\label{sec:discussion}
Three transferable lessons stand out.

\textbf{(1) Task design controls what ``memory'' means.} A single geometric
change---confining evidence to the left of a grass corridor wider than the view
radius---converts a reflex task into a working-memory task, because it forces the
evidence out of the observation window before the decision. Whenever a
partially-observed task is meant to test memory, one should check that the cue
and the response cannot co-occur in a single observation.

\textbf{(2) Discounted potential shaping can pay for loitering.} Equation
\eqref{eq:reward} with $\gamma_{\!\text{sh}}=\gamma<1$ gives a positive
per-step reward $\kappa(1-\gamma)\ctg$ for standing still, which beats a small
time cost when the potential is large. Imagination-based agents fall into the
resulting basin far more readily than on-policy ones. Using pure-progress shaping
($\gamma_{\!\text{sh}}=1$) removes the basin at no cost to the useful gradient.

\textbf{(3) Decodable, causal, and steerable are three different claims.} The
belief here is strongly decodable and strongly causal (behaviour tracks it almost
perfectly), yet only moderately steerable by linear edits, because the
representation entangles category with episode phase and curves through activation
space. Interpretability claims should say which of the three they mean.

% ======================================================================
\section{Conclusion}
\label{sec:conclusion}
Cogniland packages a simple but pointed question---\emph{what does an agent
believe about a hidden variable, and can we read, edit, and watch it plan with
that belief?}---into a small, fully-controlled navigation task and a matched set
of tools. The \code{fork\_wall} environment makes the hidden variable known, the
memory requirement real, and the decision discrete; recurrent PPO and DreamerV3
give us a policy-only and a generative agent on identical data; and linear
probing, activation steering, manifold geometry, and dream decoding let us turn
the vague notion of ``belief'' into measured, falsifiable statements. The
findings---belief drives the decision more tightly than ground truth, belief
generates the dream, belief curves through activation space but robustly---make
the environment a useful substrate for the mechanistic study of memory and
planning, and the reward-shaping and task-design pitfalls we hit and fixed are
lessons that generalize well beyond it.

% ======================================================================
\small
\begin{thebibliography}{99}
\bibitem[Cho et al.(2014)]{cho2014gru} K.~Cho, B.~van Merri\"enboer, C.~Gulcehre,
D.~Bahdanau, F.~Bougares, H.~Schwenk, and Y.~Bengio. Learning phrase
representations using RNN encoder--decoder for statistical machine translation.
\emph{EMNLP}, 2014.
\bibitem[Geiger et al.(2021)]{geiger2021causal} A.~Geiger, H.~Lu, T.~Icard, and
C.~Potts. Causal abstractions of neural networks. \emph{NeurIPS}, 2021.
(Framework realized in Goodfire's \code{causalab}, \emph{manifold\_steering}
branch.)
\bibitem[Ha \& Schmidhuber(2018)]{ha2018world} D.~Ha and J.~Schmidhuber. World
models. \emph{NeurIPS}, 2018.
\bibitem[Hafner et al.(2023)]{hafner2023dreamerv3} D.~Hafner, J.~Pasukonis,
J.~Ba, and T.~Lillicrap. Mastering diverse domains through world models.
\emph{arXiv:2301.04104}, 2023.
\bibitem[Kaelbling et al.(1998)]{kaelbling1998planning} L.~Kaelbling, M.~Littman,
and A.~Cassandra. Planning and acting in partially observable stochastic domains.
\emph{Artificial Intelligence}, 1998.
\bibitem[Lambrechts et al.(2022)]{lambrechts2022recurrent} G.~Lambrechts,
A.~Bolland, and D.~Ernst. Recurrent networks, hidden states and beliefs in
partially observable environments. \emph{TMLR}, 2022.
\bibitem[Li et al.(2023)]{li2023othello} K.~Li, A.~K. Hopkins, D.~Bau,
F.~Vi\'egas, H.~Pfister, and M.~Wattenberg. Emergent world representations:
exploring a sequence model trained on a synthetic task. \emph{ICLR}, 2023.
\bibitem[Mante et al.(2013)]{mante2013context} V.~Mante, D.~Sussillo, K.~Shenoy,
and W.~Newsome. Context-dependent computation by recurrent dynamics in prefrontal
cortex. \emph{Nature}, 2013.
\bibitem[Meng et al.(2022)]{meng2022locating} K.~Meng, D.~Bau, A.~Andonian, and
Y.~Belinkov. Locating and editing factual associations in GPT. \emph{NeurIPS},
2022.
\bibitem[Mikulik et al.(2020)]{mikulik2020meta} V.~Mikulik, G.~Del\'etang,
T.~McGrath, et al. Meta-trained agents implement Bayes-optimal agents.
\emph{NeurIPS}, 2020.
\bibitem[Morihira et al.(2026)]{morihira2026r2dreamer} N.~Morihira, A.~Nahar,
K.~Bharadwaj, Y.~Kato, A.~Hayashi, and T.~Harada. R2-Dreamer: redundancy-reduced
world models without decoders or augmentation. \emph{ICLR}, 2026.
\bibitem[Ng et al.(1999)]{ng1999policy} A.~Ng, D.~Harada, and S.~Russell. Policy
invariance under reward transformations: theory and application to reward
shaping. \emph{ICML}, 1999.
\bibitem[Ravichander et al.(2021)]{ravichander2021probing} A.~Ravichander,
Y.~Belinkov, and E.~Hovy. Probing the probing paradigm: does probing accuracy
entail task relevance? \emph{EACL}, 2021.
\bibitem[Reinsch(1967)]{reinsch1967smoothing} C.~Reinsch. Smoothing by spline
functions. \emph{Numerische Mathematik}, 1967.
\bibitem[Schulman et al.(2017)]{schulman2017proximal} J.~Schulman, F.~Wolski,
P.~Dhariwal, A.~Radford, and O.~Klimov. Proximal policy optimization algorithms.
\emph{arXiv:1707.06347}, 2017.
\bibitem[Shai et al.(2024)]{shai2024transformers} A.~Shai, S.~Marzen,
P.~Riechers, et al. Transformers represent belief state geometry in their
residual stream. \emph{arXiv}, 2024.
\bibitem[Sussillo \& Barak(2013)]{sussillo2013opening} D.~Sussillo and O.~Barak.
Opening the black box: low-dimensional dynamics in high-dimensional recurrent
neural networks. \emph{Neural Computation}, 2013.
\bibitem[Turner et al.(2023)]{turner2023activation} A.~Turner, L.~Thiergart,
et al. Activation addition: steering language models without optimization.
\emph{arXiv}, 2023.
\bibitem[Wong \& Wang(2006)]{wong2006recurrent} K.-F. Wong and X.-J. Wang. A
recurrent network mechanism of time integration in perceptual decisions.
\emph{Journal of Neuroscience}, 2006.
\bibitem[Zou et al.(2023)]{zou2023representation} A.~Zou, L.~Phan, S.~Chen, et al.
Representation engineering: a top-down approach to AI transparency. \emph{arXiv},
2023.
\end{thebibliography}

\end{document}
